\section{单圈阶示例计算}\label{sec:sample}

作为示例，本节算出对规范场两点函数有贡献的单圈图的发散部分。
由此可以看到，流时间为零处所需的参数与场重正化
恰好消去所有流时间处的奇异性。

\begin{figure}[t]
\centerline{\includegraphics[width=9.6cm]{images/diagram7}}
\caption{对 $AA$（图 $1-3$）、$LA$（图 $4$）、$LB$（图 $5,6$）
与 $LL$（图 $7$）顶点函数有贡献的单圈图。
左右两端外腿上的时间--动量--指标组合分别为
$(t,p,\mu,a)$ 与 $(s,q,\nu,b)$
（对 $A$ 腿而言没有流时间）。}
\label{fig:oneloop}
\end{figure}

\subsection{自能图的计算}\label{ssec:selfenergy}

单圈阶的自能图只有寥寥几个（见图~\ref{fig:oneloop}）。
图 $1-3$ 与流时间为零处规范场两点函数中的图相同。
因此我们仅引用其奇异部分之和的已知结果
\begin{equation}
   \left.\Gamma^{(1)}_{AA}(p)_{\mu\nu}^{ab}
   \right|_{\rm pole}
   =\delta^{ab}(\delta_{\mu\nu}p^2-p_{\mu}p_{\nu})
   {N\over16\pi^2\eps}\left({13\over6}-{1\over 2\lambda_0}\right).
   \label{eq:5.1}
\end{equation}

对于图 $4$，费曼被积函数是如下形式各项的线性组合：
\begin{equation}
   \delta^{ab}
   {Q(k,p)_{\mu\nu}\over \bigl\{k^2(k+p)^2\bigr\}^2}
   \rme^{-t\{uk^2+v(k+p)^2\}},
   \qquad u,v\in\{1,\alpha_0\},
   \label{eq:5.2}
\end{equation}
其中 $k$ 是圈动量，$Q(k,p)_{\mu\nu}$ 是 $k$ 和 $p$ 的六次齐次多项式
（动量守恒因子 $(2\pi)^D\delta(p+q)$ 以下总是略去不写）。
该图在所有 $t>0$ 及任意维数 $D$ 下都有限，
但它最终要与外传播子相乘，然后还须对 $t$ 从 $0$ 到无穷大积分。
因此积分在小 $t$ 处的行为至关重要，并可能在 $D=4$
时产生奇异性。

设 $f(t)$ 是任一光滑测试函数（例如某个外传播子），时间积分
\begin{equation}
   I(k,p)=\int_{0}^{\infty}\rmd t\, f(t)\rme^{-t\{uk^2+v(k+p)^2\}}
\label{eq:5.3}
\end{equation}
可以在大 $k$ 处展成渐近级数，其领头项为
\begin{equation}
   I(k,p)={f(0)\over(u+v)k^2}+\rmO(k^{-3}).
\label{eq:5.4}
\end{equation}
于是对流时间 $t$ 的积分改善了动量积分的发散度。
特别地，在分解式
\begin{equation}
   I(k,p)=\left\{I(k,p)-{f(0)\over(u+v)(k+r)^2}\right\}+
   {f(0)\over(u+v)(k+r)^2}
   \label{eq:5.5}
\end{equation}
（其中 $r$ 是任取的外动量）中，第一项使积分在 $D=4$ 时收敛。
剩下的是一个对数发散，因而其极点奇异性与普通费曼积分
\begin{equation}
   \delta^{ab}\delta(t){1\over u+v}\int_k
   {Q(k,0)_{\mu\nu}\over(k^2)^4(k+r)^2}
   \label{eq:5.6}
\end{equation}
的极点奇异性相同，后者用标准技术即可算出。

照此进行，便得到 $LA$ 与 $LB$ 顶点函数的发散部分：
\begin{align}
   \left.\Gamma^{(1)}_{LA}(t,p)_{\mu\nu}^{ab}
   \right|_{\rm pole}
   =g_0^2\delta^{ab}\delta(t)\delta_{\mu\nu}
   {N\over16\pi^2\eps}
   \left({3\over4}+{3\over 4\lambda_0}\right),
   \label{eq:5.7}\\[3.5pt]
   \left.\Gamma^{(1)}_{LB}(t,s,p)_{\mu\nu}^{ab}
   \right|_{\rm pole}
   =g_0^2\delta^{ab}\delta(t)\delta(s)\delta_{\mu\nu}
   {N\over16\pi^2\eps}\left(-{1\over 2\lambda_0}\right),
   \label{eq:5.8}
\end{align}
而图 $7$ 即 $LL$ 顶点函数在 $D=4$ 时有限。

\subsection{重正化}

裸耦合与规范固定参数通过
\begin{equation}
   g_0^2=\mu^{2\eps}g^2Z,
   \qquad
   \lambda_0=\lambda Z_3^{-1},
   \label{eq:5.9}
\end{equation}
与重正化参数 $g$、$\lambda$ 相联系，其中 $\mu$ 是归一化质量。
到单圈阶，重正化常数 $Z$ 与 $Z_3$ 为
\begin{align}
   Z=1-{b_0\over\eps}g^2+\rmO(g^4),
   \qquad
   b_0={N\over16\pi^2}{11\over3},
   \label{eq:5.10}\\[3pt]
   Z_3=1+{c_0\over\eps}g^2+\rmO(g^4),
   \qquad
   c_0={N\over16\pi^2}\left({13\over6}-{1\over2\lambda}\right),
   \label{eq:5.11}
\end{align}
此外还有依赖于方案的有穷项。

基本规范场的重正化
\begin{equation}
   A_{\mu}^a=Z^{1/2}Z_3^{1/2}(\Ar)^a_{\mu}
   \label{eq:5.12}
\end{equation}
同时涉及两个重正化常数，这是本文采用的非常规归一化约定所致。
出于第 \ref{sec:fin} 节将解释的原因，鬼场按非对称方式重正化：
\begin{equation}
   c^a=\tilde{Z}_3Z^{1/2}Z_3^{1/2}(\ren{c})^a,
   \qquad
   \cbar^a=Z^{1/2}Z_3^{-1/2}(\ren{\cbar})^a,
   \label{eq:5.13}
\end{equation}
其中 $\tilde{Z}_3$ 是通常的鬼场重正化常数。注意，
由于场 $c$ 与 $\cbar$ 总是成对出现、起作用的只是二者
重正化因子的乘积，式 \eqref{eq:5.13}
在流时间为零的关联函数层面等价于标准的重正化方案。

\subsection{体场需要重正化吗？}

在微扰论的单圈阶，这个问题可以通过显式计算体场关联函数的
发散部分（如果有的话）来回答。我们首先考虑 $B$ 场的两点函数，
并通过
\begin{align}
   \bigl\langle\tilde{B}_{\mu}^a(t,p)\tilde{B}_{\nu}^b(s,q)\bigr\rangle
   =(2\pi)^D\delta(p+q){\delta^{ab}\over(p^2)^2}
   \notag\\[2.5ex]
   \times
   \bigl\{(\delta_{\mu\nu}p^2-p_{\mu}p_{\nu})\mathcal{A}(t,s,p^2)
   +p_{\mu}p_{\nu}\mathcal{B}(t,s,p^2)\bigr\}
   \label{eq:5.14}
\end{align}
定义其洛伦兹不变部分 $\mathcal{A}$ 与 $\mathcal{B}$。
图~\ref{fig:oneloop} 中画出的所有自能图都对 $\mathcal{A}$ 与 $\mathcal{B}$
有贡献。图 $4-6$ 实际上各有两个贡献，因为它们的外腿不同。

在重正化微扰展开
\begin{equation}
   \mathcal{X}=\mu^{2\eps}\sum_{l=0}^{\infty}g^{2l+2}\mathcal{X}^{(l)},
   \qquad
   \mathcal{X}=\mathcal{A}\;\hbox{或}\;\mathcal{B}
   \label{eq:5.15}
\end{equation}
中，领头阶系数为
\begin{equation}
   \mathcal{A}^{(0)}=\rme^{-(t+s)p^2},
   \qquad
   \mathcal{B}^{(0)}=\lambda^{-1}\rme^{-\alpha_0(t+s)p^2}.
   \label{eq:5.16}
\end{equation}
于是在下一阶，来自耦合与规范固定参数重正化的极点的留数为
\begin{align}
   \left.\mathrm{res}\{\mathcal{A}^{(1)}\}\right|_{Z\;\mathrm{factors}}=
   -b_0\mathcal{A}^{(0)},
   \label{eq:5.17}\\[2.5pt]
   \left.\mathrm{res}\{\mathcal{B}^{(1)}\}\right|_{Z\;\mathrm{factors}}=
   (c_0-b_0)\mathcal{B}^{(0)}.
   \label{eq:5.18}
\end{align}
回顾小节 \ref{ssec:selfenergy} 得到的结果，现在容易验证这些极点
恰好被单圈图的发散部分消去。到这一阶微扰论为止，
随时间变化的规范场的两点函数是有限的，不需要进一步重正化。

对于 $BL$、$B\ren{A}$、$L\ren{A}$ 与 $LL$ 关联函数，
$D=4$ 处奇异性的消去可以用同样的方式证明。
实际上后两种情形几乎无需证明，因为这两个两点函数
按微扰论所有阶都为零（不存在只有朝内而没有朝外流线的图）。
本节报告的所有计算由此支持如下猜想：
$B$ 场与 $L$ 场不需要重正化。
